\section{Conclusion}
\label{sec:conclusion}

We introduced G-MFRM, a Bayesian measurement model that jointly estimates model ability, judge severity, and contamination susceptibility within a single inferential framework. The key theoretical contribution is formalizing the confounding between judge bias and benchmark contamination as omitted variable bias: omitting the contamination facet produces systematic bias in ability estimates, driven primarily by the specification error itself rather than the exposure-severity covariance (Section~\ref{sec:confounding}).

On Arena-Hard, G-MFRM provides interpretable per-model contamination diagnostics unavailable from existing methods: older models exhibit 4--6$\times$ higher susceptibility, and MFRM-only estimation inflates ability for the most contaminated models by up to 0.79 logits. A controlled coverage ablation ($\Delta\tau$: $-0.044 \to -0.015 \to 0.000$ as coverage increases to 100\%) demonstrates that over-correction is attributable to incomplete exposure estimation, and complete coverage eliminates the over-correction observed in partial-coverage settings. Ablation confirms joint estimation as the dominant contributor: staged estimation doubles ability RMSE.

Two practical implications emerge: evaluation frameworks that treat judge bias and contamination as independent incur systematic errors, and the coverage--degradation finding provides an actionable quality requirement for contamination-aware evaluation. For practitioners, G-MFRM can serve as a post-hoc calibration layer whose exposure pipeline (\S\ref{sec:staged}) is modular. Open directions include higher-rank contamination decompositions, simulation-based recalibration \citep{talts2018sbc} of $c_n$ credible intervals, and scaling to larger model populations.